% Part: normal-modal-logic
% Chapter: frame-correspondence
% Section: introduction

\documentclass[../../../include/open-logic-section]{subfiles}

\begin{document}

\olfileid{nml}{frd}{int}

\olsection{引言}

模态逻辑学家感兴趣的一个问题是，可达关系与具有该关系的模型中某些公式的真值之间的关系。例如，假设可达关系是自反的，即对每个 $w \in W$，都有 $Rww$。换句话说，每个世界都可从自身可达。这意味着，当 $\Box !A$ 在世界~$w$ 上为真时，$w$ 本身就在可达世界之中，因此~$!A$ 必然在这些世界上为真。因此，如果~$\mModel{M}$ 的可达关系~$R$ 是自反的，那么无论我们取哪个世界 $w$ 和哪个公式 $!A$，$\Box !A \lif !A$ 都会在该世界上为真（换言之，模式 $\Box p \lif p$ 及其所有代入实例在~$\mModel{M}$ 中都为真）。

然而，逆命题并不成立。例如，并非如果 $\Box p \lif p$ 在~$\mModel{M}$ 中为真，$R$ 就是自反的。因为我们很容易找到一个非自反的模型~$\mModel{M}$，使得 $\Box p \lif p$ 在所有世界上都为真：取一个只含单个世界~$w$ 的模型，$w$ 不可从自身可达，但 $w \in V(p)$。通过适当地选取 $p$ 的真值，我们可以在一个非自反的模型中使 $\Box !A \lif !A$ 为真。

解决方法是将赋值~$V$ 从考量中移除。如果我们要求，无论 $V(p)$ 中包含哪些世界，$\Box p \lif p$ 在~$\mModel{M}$ 的所有世界上都为真，那么 $R$ 必定是自反的。因为在任何非自反模型中，都至少存在一个世界~$w$ 使得 $Rww$ 不成立。如果我们令 $V(p) =
W \setminus \{w\}$，那么 $p$ 将在除~$w$ 以外的所有世界上为真，因此在所有从~$w$ 可达的世界上也为真（因为已知 $w$ 不可从自身可达，且 $w$ 是唯一一个~$p$ 为假的世界）。另一方面，$p$ 在 $w$ 上为假，所以 $\Box p \lif p$ 在~$w$ 上为假。

这表明，我们应当为不带赋值的模型结构引入一种记法：我们称之为\emph{框架}。框架 $\mModel{F}$ 就是由世界集和可达关系组成的二元组 $\tuple{W, R}$。于是，正如我们的说法，每个模型 $\tuple{W, R, V}$ 都\emph{基于}框架 $\tuple{W, R}$。反之，一个框架决定了基于它的所有模型构成的类；而一个框架类决定了基于该类中任意框架的所有模型构成的类。我们可以定义 $\mModel{F} \Entails !A$，即公式在框架中\emph{有效}这一概念：对所有基于~$\mModel{F}$ 的 $\mModel{M}$，都有 $\mSat{M}{!A}$。

利用这种记法，我们可以在公式与框架类之间建立对应关系：例如，$\mModel{F} \Entails \Box p \lif p$ 当且仅当 $\mModel{F}$ 是自反的。

\end{document}